\section{模型求解}

\subsection{求解流程}

当问题规模较大时，可采用启发式或智能优化算法获得近似最优解；当问题规模较小且约束较规整时，也可直接使用精确求解器。模板中保留本节，主要是为了展示算法步骤说明与伪代码环境的排版空间。

\begin{algorithm}[H]
\caption{示意性求解流程}
\KwIn{客户节点集合、车辆集合、距离矩阵与约束参数}
\KwOut{满足约束的配送方案}
初始化可行解与参数配置\;
根据目标函数计算当前方案适应度\;
循环执行邻域搜索或启发式更新，直到满足停止条件\;
输出当前最优方案并记录关键指标\;
\end{algorithm}

\subsection{实现说明}

若正文需要补充实现细节，可在这里说明求解器选择、参数搜索策略、停止准则和复杂度控制方式。比赛论文中，这一部分通常与结果分析相互呼应，用来解释为何当前算法能够在给定规模下得到稳定结果。
